\section{Thuật toán SAT tìm trạng thái Alien Tiles khó nhất}

Trong phần này, chúng tôi trình bày phương pháp giải bài toán
Alien Tiles bằng cách sử dụng bộ giải SAT kết hợp với các ràng buộc
Pseudo-Boolean.

Mục tiêu của thuật toán là tìm trạng thái mục tiêu của bảng
$c$ màu kích thước $n \times n$ sao cho số lần bấm tối thiểu
để tạo ra trạng thái đó là lớn nhất.

\subsection{Ý tưởng}

Thuật toán hoạt động theo hai bước chính:

\begin{itemize}
\item Sinh các trạng thái mục tiêu có tổng số lần bấm bằng $k$.
\item Kiểm tra xem trạng thái đó có thể được tạo ra với ít hơn $k$
lần bấm hay không.
\end{itemize}

Nếu số lần bấm tối thiểu đúng bằng $k$, trạng thái đó là hợp lệ.

Thuật toán sẽ tăng dần $k$ cho đến khi không còn trạng thái nào
thỏa mãn.

\subsection{Biến sử dụng}

\begin{itemize}
\item $x_{i,j,k}$: biến Boolean biểu diễn ô $(i,j)$ được bấm
đúng $k$ lần.
\item $s_{i,j,m}$: biến Boolean biểu diễn màu cuối cùng của ô
$(i,j)$ là $m$.
\end{itemize}

Với mỗi ô ta có ràng buộc:

\[
\sum_{k=0}^{c-1} x_{i,j,k} = 1
\]

đảm bảo mỗi ô chỉ có một số lần bấm.

Tương tự:

\[
\sum_{m=0}^{c-1} s_{i,j,m} = 1
\]

đảm bảo mỗi ô chỉ có một màu cuối cùng.

\subsection{Ràng buộc ảnh hưởng}

Khi bấm ô $(i,j)$, tất cả các ô cùng hàng và cùng cột sẽ
tăng giá trị thêm $1$ (mod $c$).

Tổng ảnh hưởng lên ô $(i,j)$ được tính như sau:

\[
S_{i,j} =
\sum_{col=1}^{n} x_{i,col}
+
\sum_{row=1}^{n} x_{row,j}
-
x_{i,j}
\]

Giá trị màu của ô phải thỏa mãn:

\[
S_{i,j} \bmod c = t_{i,j}
\]

Ràng buộc này được mã hóa bằng các ràng buộc
Pseudo-Boolean cho tất cả các giá trị tổng có thể.

\subsection{Ràng buộc tổng số lần bấm}

Tổng số lần bấm trên toàn bộ bảng được cố định:

\[
\sum_{i=1}^{n} \sum_{j=1}^{n} x_{i,j} = k
\]

trong đó $k$ là số lần bấm đang được kiểm tra.

\subsection{Thuật toán tìm trạng thái khó nhất}

\begin{algorithm}
\caption{Tìm trạng thái Alien Tiles khó nhất}
\begin{algorithmic}[1]

\State $k \gets 1$
\State $best \gets$ None

\While{còn tồn tại nghiệm}

    \State Sinh mô hình SAT với tổng số lần bấm bằng $k$

    \If{không tồn tại nghiệm}
        \State \textbf{dừng}
    \EndIf

    \State Lấy trạng thái mục tiêu $T$ từ nghiệm SAT

    \State Tính số lần bấm tối thiểu $k_{min}$ để tạo ra $T$

    \If{$k_{min} = k$}
        \State $best \gets T$
    \EndIf

    \State $k \gets k + 1$

\EndWhile

\State Trả về giá trị $k-1$

\end{algorithmic}
\end{algorithm}

\subsection{Tìm số lần bấm tối thiểu}

Với một trạng thái mục tiêu cố định $T$, số lần bấm tối thiểu
được tìm bằng cách lặp:

\begin{enumerate}

\item Giải bài toán SAT để tìm một nghiệm.

\item Tính tổng số lần bấm từ nghiệm.

\item Thêm ràng buộc:

\[
\sum x_{i,j} \le k - 1
\]

\item Giải lại bài toán.

\end{enumerate}

Quá trình lặp lại cho đến khi công thức trở nên
không thỏa mãn. Giá trị $k$ cuối cùng chính là
số lần bấm tối thiểu.

\subsection{Độ phức tạp}

Với bảng kích thước $n \times n$ và $c$ màu:

\begin{itemize}

\item Số biến bấm: $O(n^2 c)$

\item Số biến màu: $O(n^2 c)$

\item Số ràng buộc Pseudo-Boolean: $O(n^2)$

\end{itemize}

Thuật toán cần giải nhiều bài toán SAT nên
thời gian chạy phụ thuộc mạnh vào kích thước bảng
và số màu.
